% Sección: Conclusiones

Los resultados del ejercicio fueron positivos; todos los modelos alcanzaron
métricas sumamente altas en el conjunto de testeo con muy pocos casos mal
clasificados y, más importantemente en algunos casos, sin falsos negativos.
Este resultado, sin embargo, probablemente diga más sobre la calidad de los
datos que sobre el ejercicio de modelado. El conjunto de datos no contaba con
valores faltantes, no parecía haber errores en la información, y en general
había fuertes asociaciones entre las variables predictoras y la respuesta.

Uno de los dos mejores modelos es paramétrico y el otro no paramétrico. Bosques
aleatorios tuvo un desempeño marginalmente mejor, pero la regresión logística
con LASSO tiene la ventaja de que, por un lado, permite reducir la cantidad de
variables que es necesario medir para obtener una predicción y, por otro lado,
los coeficientes que estima resultan más interpretables, sobre todo en este
caso que corresponden a las variables originales normalizadas. Es posible
extraer la importancia de cada variable de un modelo de bosques aleatorios,
pero las medidas de importancia no son tan directamente interpretables como los
coeficientes de una regresión, que tienen vínculo directo con el modelo
probabilístico que usamos para conceptualizar el fenómeno. Por estos motivos
considero que es preferible el modelo logístico por sobre el de bosques
aleatorios.

Una de las principales dificultades del trabajo fue la falta de conocimiento
sustantivo sobre el área del conocimiento a la que pertenecen los datos.
Resultó difícil generar ideas a partir de las distribuciones y asociaciones
observadas en la etapa de exploración al no tener noción de la significancia
científica de las variables. Finalmente se utilizaron todas las variables
disponibles en casi todos los modelos ya que no hubo motivos sustantivos para
excluir, combinar o transformarlas más allá de lo que exigían los modelos. Esto
también impacto en la selección de la métrica de evaluación. Se optó por usar
el puntaje $F_\beta$, pero, sin conocimiento del área, resultó prácticamente
imposible decidir una ponderación adecuada. Por otra parte, el error puede
haber estado, justamente, en la elección de esa métrica.
